% !TeX program = lualatex
% !TeX encoding = utf8
% !TeX spellcheck = uk_UA
% !BIB program = biber
% !TEX indexprogram = upmendex

\documentclass[%
biblatex,
%marginversioninfo
]{ConspectBook}

%chktex-file 1

%=================================================================================
%
%							Вихідні дані докумета
%
%=================================================================================

\title{Класична електродинаміка}
\def\subtitle{Мікроскопічна теорія}


%============================ Бібліографія =======================================
%
\addbibresource{./bibliography/books.bib}
%\addbibresource{Kvant.bib}
%
%=================================================================================

%=================================================================================
%
\input{CoverPage}%			Титульна сторінка
%
%=================================================================================

\begin{document}
\pagestyle{empty}
%\ifelectronic\CoverPage\setcounter{page}{0}\fi
\CoverPage
\pagenumbering{arabic}
\maketitle
\makeinfopage
%=================================================================================
%
\clearpage\pagestyle{plain}
\tableofcontents%			     Зміст
%
%=================================================================================

%=================================================================================
%
%\include{intro}%              Вступ та передмова
%
%=================================================================================

%=================================================================================
%
%							Вставка файлів розділів
%
%=================================================================================

\addtocontents{toc}{\protect\vfill}
\input{Preface}
\input{Measures}
\pagestyle{main}
%------------------- Частина І -------------------
\part{Класична електродинаміка}
\label{part:I}
\epigraph{Неможливо позбутися відчуття, що ці математичні формули існують незалежно від нас і володіють власним розумом, що вони мудріші за нас, мудріше навіть тих, хто їх відкрив, і що ми дістаємо з них більше, ніж спочатку було закладено\ldots}{Heinrich Hertz}
\multiinclude{
	Basics,
	Solutions,
	FreeField,
	Radiation
}[]

%------------------- Частина IІ ------------------
\addtocontents{toc}{\protect\clearpage}
\part{Чотиривимірне формулювання класичної\\ електродинаміки}
\label{part:II}
\epigraph{Рівняння Максвелла
	інваріантні відносно перетворень Лоренца~--- саме тому їх
	і можна записати красиво.}{Richard Feynman,
	\textit{The Feynman Lectures on Physics}, Vol.~II, Ch.~25}
\multiinclude{
	LorentzTransform,
	Minkovsky,
	RelElectrodynamics,
	VariationPrinciple,
	TEI,
	RadiationReaction
}[]

%=================================================================================
%
%						      Вставка відповідей
%
%=================================================================================
%\bookmarksetup{startatroot}
%\answers
%\multiinclude{\ChaptersOne}[-Answers]
%\multiinclude{\ChaptersTwo}[-Answers]
%=================================================================================
%
%								Додатки
%
%=================================================================================
\appendix
\input{Additions/SGStoSI}
\input{Additions/Vanaliz}
% !TeX program = lualatex
% !TeX encoding = utf8
% !TeX spellcheck = uk_UA
% !TeX root =../ClassicalElectrodynamics.tex

% chktex-file 3
% chktex-file 21
% chktex-file 25
% chktex-file 46

%% --------------------------------------------------------
\section{Узагальнені функції}\setcounter{equation}{0}\label{sec:GeneralizedFunctions}
%% --------------------------------------------------------

%Дельта-функція Дірака (або $\delta$-функція) є узагальненою функцією і була введена фізиком Полем Діраком для моделювання густини ідеалізованої
%точкової маси або точкового заряду.
%
%На <<фізичному рівні строгості>> можна визначити $\delta$-функцію формальним співвідношенням:
%\begin{equation}
%    \int\limits_{-\infty}^{+\infty}f(x)\delta(x-x_0)\,dx=f(x_0),
%\end{equation}
%У випадку інтегрування по скінченному об'єму $V$:
%\begin{equation}\label{eq:delta3D}
%    \int\limits_{V}f(\vect{r})\delta(\vect{r} - \vect{r}_0) dV=f(\vect{r}_0),
%\end{equation}
%де точка $\vect{r}_0$ знаходиться всередині об'єму $V$.

%$\delta$-функцію однієї дійсної змінної можна визначити як функцію, що задовольняє наступним умовам:
%\begin{equation}\label{eq:delta}
%\delta(x - x_0)=\left\{\begin{matrix}
%   +\infty, & x=x_0, \\
%   0, & x \neq x_0; \\
%\end{matrix}\right.
%\end{equation}
%\begin{equation}\label{eq:deltaProp0}
%    \int\limits_{-\infty}^{+\infty}\delta(x-x_0)dx=1.
%\end{equation}
%Тобто ця функція не дорівнює нулю тільки в точці $x = x_0$, де вона перетворюється в нескінченність таким чином, щоб її інтеграл в будь-якому околі точки $x = x_0$ дорівнює $1$.
%
%Аналітичне представлення $\delta$-функції:
%\begin{equation}\label{eq:deltaanalit}
%    \delta(x) = \frac{1}{2\pi}\int\limits_{-\infty}^{\infty} e^{ikx}dk.
%\end{equation}
%\bigskip\noindent%
%\textbf{Властивості дельта-функції}
%\bigskip
%
%\begin{enumerate}[label=\alph*)]
%\item Дельта-функція парна $\delta(-x) = \delta(x)$,
%\item $x\delta(x) = 0$,
%\item $\delta(ax) = \frac{1}{|a|}\delta(x)$,
%\item  $x\delta^\prime(x)=-\delta(x)$,
%\item $\delta(f(x))=\sum\limits_k\frac{\delta(x-x_k)}{|f'(x_k)|}$, де $x_k$~--- нулі функції $f(x)$,
%
%\end{enumerate}



В задачах електродинаміки часто виникає необхідність розглядати заряджені тіла, розміри яких дуже малі у порівнянні з іншими просторовими масштабами. З
цим пов'язана модель точкового (сферичного) заряду, густину якого подають у вигляді:

\begin{equation}
\rho(\vect{r}) = q \delta(\vect{r} - \vect{r}_0),
\end{equation}
%
де \( q \) --- величина заряду, \( \vect{r} \) --- його положення, \( \delta \) --- функція Дірака, що визначається умовою:

\begin{equation}
\int dV \, \delta(\vect{r} - \vect{r}_0) \chi(\vect{r}) = \chi(\vect{r}_0)
\label{eq:dirac_definition}
\end{equation}
%
для будь-якої неперервної функції \( \chi \).

Очевидно, що визначення~\eqref{eq:dirac_definition} \(\delta\)-функції не є математично коректним, якщо вважати, що в~\eqref{eq:dirac_definition} маємо
справу зі звичайним інтегруванням, наприклад, за Лебегом. Цьому співвідношенню можна надати математичний зміст, якщо розглядати \(\delta\)-функцію як
слабку границю деякої послідовності звичайних функцій. Якщо покласти (в одновимірному варіанті):

\begin{equation}
\delta_{\varepsilon}(x) = \frac{1}{\varepsilon \sqrt{\pi}} e^{-x^2 / \varepsilon^2},
\label{eq:delta_sequence}
\end{equation}

то границя:

\begin{equation}
\lim_{\varepsilon \to 0} \int dx \, \delta_{\varepsilon}(x) \chi(x) = \chi(0)
\label{eq:delta_limit}
\end{equation}
існує для будь-якої послідовності \( \varepsilon_n \to 0 \), \( n \to \infty \) та будь-якої функції \( \chi(x) \), що задовольняє досить широким умовам
(наприклад, якщо \( \chi(x) \) --- гладка та обмежена). На практиці зручніше розглядати подібні границі за жорсткіших обмежень на \( \chi(x) \).

Таке визначення цілком відповідає фізичним уявленням про точковий заряд, зосереджений в нескінченно малій області. Насправді для багатьох застосувань
істотно лише те, що розміри заряду значно менші за відстань до нього. Разом із тим треба пам'ятати, що ці міркування не проходять, коли треба розглядати
співвідношення, нелінійні за густиною заряду або за напруженостями полів; наприклад, у формулах для енергії.

%% --------------------------------------------------------
\subsection{Узагальнені та основні функції (\texorpdfstring{\( x \in \mathbb{R}^n \)}{x ∈ Rⁿ})}
%% --------------------------------------------------------

Узагальнену \(\delta\)-функцію можна визначити за допомогою іншої, відмінної від~\eqref{eq:delta_sequence}, послідовності, тобто означення
\eqref{eq:delta_sequence}, \eqref{eq:delta_limit} не є єдиним (див. приклади далі). Але усі подібні співвідношення визначають лінійний неперервний
функціонал, що зіставляє функції \( \chi(x) \) число \( \chi(0) \). Далі розглянемо лінійні неперервні функціонали на деякій множині функцій, які будемо
називати \textit{основними}.

Нехай множина основних функцій --- це множина \( \mathbb{D} \) фінітних (тобто відмінних від нуля у деяких обмежених областях) нескінченно диференційовних
функцій. Збіжність в \( \mathbb{D} \) визначають так: послідовність \( \{\phi_k\} \subset \mathbb{D} \) збігається до \( \phi \in \mathbb{D} \), якщо всі \( \phi_k \) можуть бути
відмінними від нуля лише у деякій спільній обмеженій області та усі похідні рівномірно збігаються до відповідних похідних від \( \phi \).

Узагальнена функція --- це лінійний неперервний функціонал, означений на \( \mathbb{D} \). Множину таких узагальнених функцій позначають через \( D' \).

Існує багато способів регуляризації дельта-функції Дірака; зокрема, використовуються наступні наближення:

\begin{equation}
    \delta_{\varepsilon}(x) = \frac{1}{\pi} \frac{\sin(x/\varepsilon)}{x},
\label{eq:delta_sin}
\end{equation}

\begin{equation}
    \delta_{\varepsilon}(x) = \frac{1}{\pi} \frac{\varepsilon}{x^2 + \varepsilon^2},
\label{eq:delta_lorentz}
\end{equation}

\begin{equation}
    \delta_{\varepsilon}(x) = \frac{1}{\pi} \frac{\sin^2(x/\varepsilon)}{x^2},
\label{eq:delta_sin2}
\end{equation}
де \( x \in \mathbb{R} \).

Далі будемо позначати значення функціоналу, що відповідає узагальненій функції \( F \in D' \) на основній функції \( \chi \in \mathbb{D} \), як \( (F, \chi) \).
Наприклад, співвідношення \( (\delta, \chi) = \chi(0) \) визначає \(\delta\)-функцію Дірака. Узагальнену функцію \( F(x) \) (\( x \in \mathbb{R}^n \))
називають \textit{регулярною}, якщо існує інтегровна функція \( f(x) \), така, що відповідний функціонал можна подати у вигляді звичайного інтеграла
Лебега:

\begin{equation}\label{eq:regular_functional}
    (F, \chi) = \int dx \, f(x) \chi(x).
\end{equation}

Якщо це неможливо, \( F \) називають \textit{сингулярною узагальненою функцією}. Прикладом такої функції є \(\delta\)-функція Дірака. Проте, у
фізичній літературі для сингулярних функцій також використовують запис~\eqref{eq:regular_functional}, маючи на увазі певний граничний перехід --- типу
\eqref{eq:delta_limit} або інший. Можна показати, що будь-яку узагальнену функцію можна подати як слабку границю послідовності основних функцій з~\( D
\).

%Далі ми також будемо застосовувати формальний запис~\eqref{eq:regular_functional}, пам'ятаючи про вказані застереження.

Визначимо похідну \( \partial_i F(x) \) від узагальненої функції \( F \) співвідношенням:

\begin{equation}\label{eq:derivative_definition}
    (\partial_i F, \chi) = - (F, \partial_i \chi), \quad x \in \mathbb{R}^n.
\end{equation}

Для регулярної диференційовної функції \( F \) це співвідношення є очевидним наслідком інтегрування за частинами:

\begin{equation*}
    \int dx \, \partial_i F(x) \chi(x) = - \int dx \, F(x) \partial_i \chi(x),
\end{equation*}
де межові члени зникають, оскільки \( \chi(x) \) є фінітною.


%% --------------------------------------------------------
\section{Фундаментальні розв'язки}
%% --------------------------------------------------------


Нехай \( \hat{L} \) --- диференціальний оператор скінченого порядку зі сталими коефіцієнтами, \( x \in \mathbb{R}^n \).

\emph{Фундаментальним розв'язком оператора \( \hat{L} \)} називають узагальнену функцію \( G \), таку, що:

\begin{equation}\label{eq:fundamental_solution}
    \hat{L} G = \delta(x).
\end{equation}

%% --------------------------------------------------------
\subsection{Фундаментальний розв'язок оператора Лапласа}
%% --------------------------------------------------------

Для оператора Лапласа в \( \mathbb{R}^3 \) (\( x = \vect{r} = (x_1, x_2, x_3) \), \( r = |\vect{r}| \)) фундаментальний розв'язок має вигляд:

\begin{equation}\label{eq:laplace_G}
    G(\vect{r}) = -\frac{1}{4\pi r}, \qquad \Laplasian G = \delta(\vect{r}),
\end{equation}

або, що еквівалентно:

\begin{equation}\label{eq:laplace_fundamental}
    \Laplasian \left( \frac{1}{r} \right) = -4\pi \delta(\vect{r}).
\end{equation}

Це співвідношення перевіряють, застосувавши до $1/r$ означення похідної~\eqref{eq:derivative_definition} та виконавши інтегрування за частинами у сферичних координатах; кутові внески дають нуль, а радіальний інтеграл --- межовий доданок при $r=0$, що і дає $-4\pi\chi(0)$.

%% --------------------------------------------------------
\subsection{Фундаментальний розв'язок оператора Даламбера (хвильового оператора)}
%% --------------------------------------------------------

Нехай \( c = 1 \), \( \Dalambertian = \dfrac{\partial^2}{\partial t^2} - \Laplasian \), \( x = (t, \vect{r}) \in \mathbb{R}^4 \), \( r = |\vect{r}| \).
Фундаментальний розв'язок (запізнювальна функція Гріна) має вигляд:

\begin{equation}\label{eq:dalambert_fundamental}
    G(x) = \frac{1}{4\pi r} \delta(t - r), \qquad \Dalambertian G = \delta(x).
\end{equation}

Перевірка зводиться до підстановки $G$ у~\eqref{eq:derivative_definition}, переходу до сферичних координат і інтегрування по $t$ за допомогою $\delta(t-r)$: кутова частина оператора Даламбера дає нуль після інтегрування по сфері, а радіальна частина зводиться до повної похідної, межові значення якої дають $\chi(0,\vect{0})$.

%% --------------------------------------------------------
\subsection{Фундаментальний розв'язок оператора Гельмгольца}
%% --------------------------------------------------------

Для оператора Гельмгольца \( \hat{L} = \Laplasian + k^2 \) (\( k \in \mathbb{R} \), \( k \neq 0 \)) фундаментальним розв'язком (запізнювальною функцією Гріна) є:

\begin{equation}\label{eq:helmholtz_fundamental}
    G(\vect{r}) = -\frac{e^{ikr}}{4\pi r}, \qquad (\Laplasian + k^2) G = \delta(\vect{r}), \qquad r = |\vect{r}|.
\end{equation}

Перевірка виконується аналогічно до випадку оператора Лапласа: при $r>0$ функція $e^{ikr}/r$ задовольняє однорідне рівняння Гельмгольца, тому весь внесок зосереджується в нескінченно малому околі $r=0$, де $e^{ikr}/r \to 1/r$, і відтворюється результат~\eqref{eq:laplace_fundamental}. При $k\to 0$ вираз~\eqref{eq:helmholtz_fundamental} переходить у~\eqref{eq:laplace_G}.

%% --------------------------------------------------------
\subsection{Згортка}
%% --------------------------------------------------------

Знайдений фундаментальний розв'язок $G$ сам по собі відповідає точковому джерелу $\delta(x)$. Щоб отримати розв'язок рівняння $\hat{L}\phi=f$ для довільної (досить регулярної) правої частини $f$, потрібно \enquote{просумувати} внески від точкових джерел, розподілених із густиною $f$ --- саме цю операцію і виконує згортка.

Згортка \( \zeta = F * \chi \) узагальненої функції \( F \) з основною \( \chi \) є за визначенням:

\begin{equation}\label{eq:convolution}
\zeta(y) = \int dx \, F(x) \chi(y - x).
\end{equation}

За допомогою згортки та фундаментального розв'язку можна будувати розв'язки рівнянь виду:

\begin{equation}
\hat{L} \phi = f, \quad f \in \mathbb{D}.
\end{equation}

Дійсно, якщо покласти \( \phi = G * f \), де \( \hat{L} G = \delta(x) \), маємо:

\begin{equation*}
\hat{L}\phi(y) = \int dx \, G(x) \hat{L}_y f(y - x) = \int dx \, \hat{L}_x G(x) f(y - x) = (\delta(x), f(y - x)) = f(y).
\end{equation*}
Тут введено індекси, щоб підкреслити, що \( \hat{L}_x \) діє на змінну \( x \), \( \hat{L}_y \) --- на \( y \). Перехід \( \hat{L}_y f(y - x) = \hat{L}_x f(y - x) \) використовує те, що \( \hat{L} \) має сталі коефіцієнти: похідна по \( y_i \) від \( f(y-x) \) збігається з похідною по \( x_i \). Цей результат за певних умов можна розширити на випадок \( f \in \mathbb{D}' \).

Звідси маємо такі розв'язки рівнянь:

\begin{equation}
\Laplasian \phi = -4\pi \rho: \phi(\vect{r}) = \int dV' \, \frac{\rho(\vect{r}')}{|\vect{r} - \vect{r}'|},
\end{equation}

\begin{equation}
\Dalambertian \phi = 4\pi \rho: \phi(t, \vect{r}) = \int \frac{\rho(t - |\vect{r} - \vect{r}'|, \vect{r}')}{|\vect{r} - \vect{r}'|} dV' ,
\end{equation}

\begin{equation}
\Laplasian \phi + k^2 \phi = -4\pi \rho: \phi(\vect{r}) = \int \frac{e^{ik|\vect{r} - \vect{r}'|}}{|\vect{r} - \vect{r}'|}
\rho(\vect{r}')\,
dV'.
\end{equation}
% !TeX program = lualatex
% !TeX encoding = utf8
% !TeX spellcheck = uk_UA
% !TeX root =../ClassicalElectrodynamics.tex

% chktex-file 3
% chktex-file 21
% chktex-file 25
% chktex-file 46

%% --------------------------------------------------------
\section{Спеціальні функції}\label{sec:special_functions}
%% --------------------------------------------------------

%% --------------------------------------------------------
\subsection{Поліноми Лежандра}\label{Polinoms}
%% --------------------------------------------------------

Поліноми Лежандра застосовуються у теорії потенціалу при розкладанні виразу в околі точки $\vect{r}$:
\begin{equation*}
    \frac{1}{|\vect{r} - \vect{r}_0|} = \frac{1}{\sqrt{r^2 - 2rr_0 \cos\chi + r_0^2 }}  =  \sum\limits_{l = 0}^{\infty} \frac{r^l_<}{r^{l+1}_>} P_l(\cos\chi),
\end{equation*}
де $r_>$, $r_<$~--- більша і менша із величин $|\vect{r}|$ та $\vect{r}_0$, відповідно, $\cos\chi$~--- кут між векторами $\vect{r}$ та $|\vect{r}_0|$.

\medskip%
\textbf{Деякі поліноми Лежандра}

\begin{align*}
P_{0}(\cos\chi)  = 1, &\quad P_{1}(\cos\chi)  = \cos\chi, \\
P_{2}(\cos\chi)  = \frac {1}{2}(3\cos^2\chi-1), &\quad P_{3}(\cos\chi)  = \frac {1}{2}(5\cos^2\chi-3\cos\chi).
\end{align*}

%% --------------------------------------------------------
\subsection{Сферичні гармоніки}\label{Spherical_Harmonics}
%% --------------------------------------------------------

Сферичні функції, що залежать від полярних кутів визначаються формулою:
\begin{equation*}
    Y_{lm}(\theta,\phi) = (-1)^{(m+|m|)/2} \sqrt{\frac{2l+1}{4\pi}\frac{(l-m)!}{(l+m)!}}P_l^{|m|}(\cos \theta) e^{im\phi},
\end{equation*}
де $l = 0,1,2, \ldots$, $m$ пробігає значення від $-l$ до $l$, а $P_l^{|m|}(x)$~--- приєднані функції Лежандра.

Вони утворюють повну ортонормовану систему функцій:
\[
    \int\limits_{\theta=0}^{\pi} \int\limits_{\phi=0}^{2\pi} Y^*_{l',m'}(\theta, \phi) Y_{l,m}(\theta, \phi) \sin\theta d\theta d\phi  = \delta_{l,l'} \delta_{m,m'}.
\]

Деякі сферичні гармоніки:
\begin{align*}
	Y_{0,0}(\theta,\phi)     & =\sqrt{\frac{1} {4  \pi}},                                            \\
	Y_{1,\pm 1}(\theta,\phi) & =\sqrt{\frac{3} {8  \pi}} \, \sin\theta \, e^{\pm i\phi},             \\
	Y_{1,0}(\theta,\phi)     & =\sqrt{\frac{3} {4  \pi}}\, \cos\theta ,                              \\
	Y_{2,0}(\theta,\phi)     & =\sqrt{\frac{5} {16 \pi}}\, (3\cos^{2}\theta-1),                     \\
	Y_{2,\pm 1}(\theta,\phi) & =\sqrt{\frac{15}{8  \pi}}\, \sin\theta\, \cos\theta\, e^{\pm i\phi}, \\
	Y_{2,\pm 2}(\theta,\phi) & =\sqrt{\frac{15}{32 \pi}} \, \sin^{2}\theta \, e^{\pm2i\phi}.
\end{align*}

%% --------------------------------------------------------
\subsection{Циліндричні функції}\setcounter{equation}{0}
%% --------------------------------------------------------

Рівняння, що виникають в задачах з циліндричною симетрією, мають вигляд:
	\begin{equation}\label{eq:Bessel_eq}
		\frac{d^2y}{dx^2} + \frac1x\frac{dy}{dx} + \left(1 - \frac{m^2}{x^2} \right) y = 0,
	\end{equation}
	розв'язок яких можна представити за допомогою функцій Бесселя $J_m(x)$ та Неймана $N_m$ у вигляді лінійної комбінації $y(x) = A J_m(x) + B N_m(x)$ або у вигляді лінійної комбінації $y(x) = A H^{(1)}_m(x) + B H^{(2)}_m(x)$, де функції $H_m^{(1,2)} =J_m \pm i N_m$~--- називаються функціями Ганкеля 1-го та 2-го роду, відповідно. Доцільність введення функцій Ганкеля обумовлена тим, що вони мають прості асимптотичні розкладання при $|x| \gg 1$ і зручні для задач, пов'язаних з поширенням хвиль.

Для $m = 0,1,2,\ldots$ функції Неймана нескінченні в точці $x = 0$, тобто $\lim\limits_{x\to0}N_m(x) = -\infty$.

Функції Бесселя можна представити за допомогою ряду (в околі точки $x = 0$ для цілих, або невід'ємних $m$):
\begin{equation}\label{eq:J}
    J_m(x) = \sum_{n=0}^\infty \frac{(-1)^n}{n! \Gamma(n+m+1)} {\left({\frac{x}{2}}\right)}^{2n+m},
\end{equation}
де $ \Gamma$~--- \href{https://en.wikipedia.org/wiki/Gamma_function}{гамма-функція}. Для $m \in \mathbb{Z}$ повинна виконуватись рівність $J_{-m}(x)  = (-1)^m J_m(x)$.

\begin{center}
  \begin{tikzpicture}
    \begin{axis}[axis lines = middle,
			axis line style={-stealth},
			minor grid style = {line width=.1pt,draw=gray!10},
            width=\textwidth, height=0.5*\textwidth, xlabel=$x$,
            xtick=\empty,
%            ytick={-0.5,0.5,1},
            legend style={draw=none},
    ]
    \addplot+[id=parable,domain=-20:20, samples=500, mark=none, width=2pt, color=red, thick]
    gnuplot{besj0(x)};% node[pin=95:{$J_0(x)$}]{};
    \addplot+[id=parable,domain=-20:20, samples=500, mark=none, width=2pt, color=blue, thick]
    gnuplot{besj1(x)};% node[pin=130:{$J_1(x)$}]{};
    \addplot+[id=parable2,domain=-20:20, samples=500, mark=none, width=2pt, color=green!50!black, thick]
    gnuplot{2*1/x*besj1(x)-besj0(x)};% node[pin=-140:{$J_2(x)$}]{};
    \legend{$J_0(x)$,$J_1(x)$,$J_2(x)$}
   \end{axis}
  \end{tikzpicture}

    {Графіки функцій Бесселя $J_m$ для $m = 0,1,2$.}
\end{center}

\begin{center}
  \begin{tikzpicture}
    \begin{axis}[axis lines = middle,
			axis line style={-stealth},
			minor grid style = {line width=.1pt,draw=gray!10},
            width=\textwidth, height=0.5*\textwidth, xlabel=$x$,
            xtick=\empty,
%            ytick={-0.5,0.5,1},
            legend style={at={(current axis.south east)}, anchor = south east, draw=none},
    ]
    \addplot+[id=parable,domain=0:40, restrict y to domain=-1.5:1, samples=500, mark=none, width=2pt, color=red, thick]
    gnuplot{besy0(x)};% node[pin=95:{$J_0(x)$}]{};
    \addplot+[id=parable,domain=0:40, restrict y to domain=-1.5:1, samples=1000, mark=none, width=2pt, color=blue, thick]
    gnuplot{besy1(x)};% node[pin=130:{$J_1(x)$}]{};
    \addplot+[id=parable2,domain=0:40, restrict y to domain=-1.5:1, samples=1000, mark=none, width=2pt, color=green!50!black, thick]
    gnuplot{2*1/x*besy1(x)-besy0(x)};% node[pin=-140:{$J_2(x)$}]{};
    \legend{$N_0(x)$,$N_1(x)$,$N_2(x)$}
   \end{axis}
  \end{tikzpicture}

    {Графіки функцій Неймана $N_m$ для $m = 0,1,2$.}
\end{center}

Функції Неймана визначаються через функції Бесселя як:
\begin{equation}\label{eq:NG}
    N_m(x) = \frac{J_m(x)\cos m\pi - J_{-m}(x)}{\sin m\pi}.
\end{equation}


Деякі рекурентні співвідношення:
\begin{equation}
    J_{m + 1}(x) + J_{m - 1}(x) = \frac{2m}{x}J_m(x).
\end{equation}

Деякі диференціальні та інтегральні співвідношення для нецілих $m$ (для цілих $m$ ці функції можна визначити за допомогою граничного переходу).:
\begin{align}
    \frac{d}{dx} J_0(x) = - J_1(x), \\
    \frac{d}{dx} \left( x^{-m}J_m(x)\right)  = - x^{-m}J_{m+1}(x), \\
     \int\limits_0^{x} x^{\prime m+1} J_m(x') dx' &= x^{m+1}J_{m+1}. \label{eq:recInt}
\end{align}

Інтеграли від добутків:
\begin{equation}
    \int\limits_0^x  J_m(k_1x')J_m(k_2x') x' dx' = \frac{x\left( k_2J_m(k_1x)J'_m(k_2x) - k_1J_m(k_2x)J'_m(k_1x)\right)}{k_1^2-k_2^{2}}   \label{eq:JJ0*}.
\end{equation}

В задачах, зазвичай, часто необхідно знайти наближений вигляд циліндричних функцій при малих та великих значеннях аргументу $x$:

при $|x| \ll 1$ з~\eqref{eq:J}

\begin{equation}
    J_0(x) \approx 1 - \frac{x^2}{4}, \quad
    J_m \approx \frac{x^m}{2^m m!}, \ m \ge 1,\, m \in \mathbb{N};
\end{equation}

при $|x| \gg 1$

\begin{align}
    J_m &\approx \sqrt{\frac{2}{\pi x}}  \cos\left( x - m\frac{\pi}{2} - \frac{\pi}{4}\right),
      \label{eq:Jxgg1}\\
    N_m &\approx \sqrt{\frac{2}{\pi x}}  \sin\left( x - m\frac{\pi}{2} - \frac{\pi}{4}\right),
  \label{eq:Yxgg1}\\
    H_m^{(1,2)} &\approx \sqrt{\frac{2}{\pi x}}  e^{\pm i \left( x - m\frac{\pi}{2} - \frac{\pi}{4}\right) }. \label{eq:Hxgg1}
\end{align}

Співвідношення Якобі-Ангера (розкладання за функціями Бесселя):
\begin{equation}\label{eq:JacobiAnger}
    e^{ix\cos\theta} = \sum\limits_{m= - \infty}^{\infty} i^mJ_m(x)e^{im\theta}, \quad
    e^{ix\sin\theta} = \sum\limits_{m= - \infty}^{\infty} J_m(x)e^{im\theta}. \\
\end{equation}
%\input{Additions/PhysConstants}
%
%%================================================================================
%
%						Вставка бібліографії
%
%=================================================================================
\clearpage\pagestyle{bibliography}

\nocite{*}

%\printbibheading[]
%
%\defbibheading{subheading}[\bibname]{%
%	\section*{#1\label{BibPage}}
%}

\printbibliography[heading=bibintoc]

\printindex

%\makelastpage
\end{document}
